\documentclass[11pt]{article}

\usepackage[margin=1in]{geometry}
\usepackage[T1]{fontenc}
\usepackage[utf8]{inputenc}
\usepackage{lmodern}
\usepackage{booktabs}
\usepackage{array}
\usepackage{longtable}
\usepackage{hyperref}
\usepackage{enumitem}

\title{Toward Behavioral Forensics of AIBO ERS-7 MIND 2:\\
Recovery, Preservation, and Hypothesis-Driven Modeling of Persistent State}
\author{Draft Manuscript}
\date{July 1, 2026}

\begin{document}
\maketitle

\begin{abstract}
This paper reports a practical research program for studying higher-order
behavior in Sony AIBO ERS-7 MIND 2 systems from preserved Memory Stick state.
The work begins with forensic recovery of a non-booting Sony-formatted stick,
establishes a reproducible restoration and preservation workflow, and then
pivots into behavior forensics over persistent application state. We describe a
host-side MIND 2 behavior simulator used as a hypothesis engine, identify
candidate persistent-state levers including \texttt{STTLOG} row \texttt{0400},
\texttt{STTLOG} row \texttt{1200}, and specimen-side \texttt{IEG.CFG}, and
show that these levers support distinct behavioral interpretations under a
growing family of probe scenarios. The current evidence suggests that
\texttt{0400} acts as a shutdown-side pivot, \texttt{1200} expresses
socially-cued persistence, and specimen-side \texttt{IEG.CFG} expresses a more
routine-preserving and multi-cycle persistence family. The contribution is not
a claim of decoded Sony internal logic, but a reproducible methodology that
links filesystem forensics, preserved media, host-side simulation, and live
hardware testing.
\end{abstract}

\section{Introduction}
Behavior in legacy autonomous robots is often treated as opaque runtime magic,
especially when vendor tooling and source-level documentation are incomplete or
when deployed systems survive only as preserved media. For the Sony AIBO ERS-7
MIND 2 platform, this creates a two-part challenge. First, one must preserve
and recover genuine bootable media accurately. Second, one must determine
whether higher-order behavioral differences are represented in persistent state
strongly enough to be isolated, transferred, and studied.

This paper presents the current state of an open research effort addressing
both problems in sequence. The first milestone established that a Sony-formatted
32 MB class stick associated with an ERS-7 startup failure could be repaired by
restoring a valid MIND 2 payload while preserving the underlying Sony media
format. The second phase treated the repaired, working image as a behavioral
specimen and developed a hypothesis-driven analysis loop over persistent files,
transplanted state trees, and host-side simulated probe scenarios.

\section{Research Goals}
The work reported here is driven by three goals:
\begin{enumerate}[leftmargin=*]
\item preserve accurate knowledge of genuine ERS-7 MIND 2 media and recovery;
\item identify which persistent files may encode behaviorally meaningful state;
\item build a reproducible path from preserved state to live-hardware
interpretation.
\end{enumerate}

The central research question is:

\begin{quote}
Where do higher-order MIND 2 behaviors live in persistent stick state, and to
what extent are those behaviors reproducible?
\end{quote}

\section{Recovered Specimen and Preservation Milestone}
The first major result was the repair and preservation of a Sony-formatted MIND
2 stick associated with a live ERS-7 startup failure. The forensic session
established four points:
\begin{enumerate}[leftmargin=*]
\item the captured medium had correct Sony 32 MB class low-level formatting;
\item the captured stick was non-bootable because it lacked
\texttt{MEMSTICK.IND} and a complete \texttt{OPEN-R/} payload tree;
\item restoring the MIND 2 payload onto that Sony-formatted medium produced a
stick that booted successfully on real hardware;
\item the repaired working stick was preserved as a new raw image for later
study.
\end{enumerate}

This matters methodologically because it separates media-format questions from
payload-state questions. The subsequent behavior work is grounded in a specimen
that is both preserved and known to boot.

\section{Methodology}
\subsection{Behavior Forensics Workflow}
The current workflow follows an observe--map--test loop:
\begin{enumerate}[leftmargin=*]
\item observe a live behavior carefully and record it with a stable template;
\item map candidate behavioral differences to persistent files and file groups;
\item create one-variable test trees and compare their outcomes.
\end{enumerate}

The preserved specimen image is held constant as the primary reference point.
Candidate state changes are staged into sacrificial trees or test media rather
than altering the preserved baseline.

\subsection{Host-Side Simulator}
Because direct robot testing is slow and expensive, a host-side MIND 2
behavior simulator was built as a hypothesis engine rather than a retail
binary emulator. The simulator reads preserved persistent state, derives a
symbolic behavior profile, and replays probe scenarios such as boot,
interaction, disruption, shutdown, and sleep requests.

The simulator is explicitly limited:
\begin{itemize}[leftmargin=*]
\item it does not execute Sony retail binaries;
\item it is not evidence of exact Sony internal implementation;
\item it is used to expose candidate structure and rank live-hardware tests.
\end{itemize}

\subsection{WLAN Service Boundary}
The current repo and recovered media also establish a useful boundary for
network-facing work. On compatible ERS-7 WLAN configurations, the robot can
expose live services including a wireless console in \texttt{WCONSOLE}, stock
MIND 2 web content, and sample HTTP camera serving. In addition, decoded
\texttt{MMFTP.CFG} content suggests bounded network-mediated content channels
for specific payload families.

Methodologically, this is important because it distinguishes two different
claims:
\begin{enumerate}[leftmargin=*]
\item the running robot can expose observable and potentially interactive WLAN
services;
\item the full Sony Memory Stick payload can be rewritten freely in place over
WLAN.
\end{enumerate}

The first claim is supported by current evidence. The second is not yet
established. Accordingly, this work treats WLAN primarily as an observation,
control, and service-interaction path rather than a proven general-purpose
remote programming path for the retail stick.

\subsection{Candidate Persistent-State Levers}
Current analysis has focused on:
\begin{itemize}[leftmargin=*]
\item \texttt{STTLOG} rows, especially \texttt{0400} and \texttt{1200};
\item \texttt{IEG.CFG} as a structured behavioral state file;
\item supporting state files including \texttt{PAT.LOG}, \texttt{FVAR},
\texttt{GVAR}, \texttt{AWAKING.CFG}, and \texttt{SIDRDATA.BIN}.
\end{itemize}

\section{Current Findings}
\subsection{\texttt{0400} as a Shutdown Pivot}
The strongest current conclusion about \texttt{STTLOG} is that row
\texttt{0400} acts as the dominant shutdown-side pivot in the host-side model.
It does not explain all behavior by itself, but it is the clearest single
lever for differentiating immediate compliance from resistance-oriented state.

\subsection{\texttt{1200} as a Social Persistence Signal}
Row \texttt{1200} currently behaves less like a routine-preserving lever and
more like a socially-cued bridge signal. Under probe scenarios involving
high-salience interaction, it expresses short-lived but meaningful persistence.
With later tuning, the model was able to sustain two cycles of social
redirection after a strong cue while remaining weaker under lower-salience
interaction.

\subsection{Specimen-Side \texttt{IEG.CFG} as a Routine-Preserving Signal}
\texttt{IEG.CFG} emerged as the strongest non-\texttt{STTLOG} behavioral
candidate. The current specimen-side file differs sharply from the bundled
baseline despite sharing the same size and \texttt{IEPR} header. When the body
after offset \texttt{0x40} is treated as a structured grid, the baseline file
has mixed non-zero values while the specimen-side file is dominated by
repeated \texttt{48.0 + 0} cells and contains a long contiguous zero-second run.

This provided the structural basis for modeling \texttt{IEG.CFG} not as a
simple boolean presence, but as a routine-rigidity signal whose strength can be
varied by partial restoration.

\section{From Hard Classifier to Gradient Model}
The first \texttt{IEG.CFG} heuristic used a hard classifier:
\begin{itemize}[leftmargin=*]
\item high zero-word density in the body;
\item high frequency of \texttt{0x42480000} (\texttt{48.0}) words;
\item resulting classification: \texttt{sparse-fixed-schedule} or not.
\end{itemize}

This was useful as a first pass and yielded a measurable threshold. In the
current restoration matrix, routine-preserving behavior remained intact through
\texttt{61} restored cells and collapsed at \texttt{62} under the old model.

The next step replaced the pure classifier with a gradient:
\begin{itemize}[leftmargin=*]
\item a continuous \texttt{ieg\_rigidity\_gradient};
\item a derived \texttt{routine\_rigidity\_signal};
\item a runtime \texttt{routine\_hold} that decays over time and disruption.
\end{itemize}

This change preserved the broad distinction between the specimen routine path
and the non-routine paths, but created a more realistic transitional band in
which partial restorations retain reduced routine influence instead of falling
instantly to zero.

\section{Dynamic Probe Results}
The most informative current probe family is time-based rather than
single-event. Three probe types are especially important:
\begin{itemize}[leftmargin=*]
\item routine endurance under repeated idle and disruption pressure;
\item social endurance after a high-salience interaction;
\item low-salience endurance after weaker interaction.
\end{itemize}

Table~\ref{tab:dynamic} summarizes the current dynamic conclusion boundary.

\begin{table}[ht]
\centering
\begin{tabular}{>{\raggedright\arraybackslash}p{0.23\linewidth} >{\raggedright\arraybackslash}p{0.18\linewidth} >{\raggedright\arraybackslash}p{0.18\linewidth} >{\raggedright\arraybackslash}p{0.18\linewidth}}
\toprule
Probe family & \texttt{0400} & \texttt{0400 + 1200} & \texttt{0400 + IEG} \\
\midrule
Routine endurance & breaks, sleeps & breaks, sleeps & preserves both cycles, stays awake \\
High-salience social endurance & breaks, sleeps & two social redirections, then sleep & preserves routine, stays awake \\
Low-salience social endurance & breaks, sleeps & weakens sooner, then sleep & preserves routine, stays awake \\
\bottomrule
\end{tabular}
\caption{Current dynamic boundary among the three main persistent-state paths.}
\label{tab:dynamic}
\end{table}

The current interpretation is:
\begin{enumerate}[leftmargin=*]
\item \texttt{0400} is a shutdown-oriented pivot;
\item \texttt{1200} supports a socially-cued, salience-sensitive persistence
family;
\item specimen-side \texttt{IEG.CFG} supports a routine-preserving,
multi-cycle persistence family.
\end{enumerate}

\section{Live-Hardware Test Matrix}
The research has now progressed far enough to justify a disciplined live bench
matrix. Four primary stick trees have been prepared:
\begin{enumerate}[leftmargin=*]
\item baseline;
\item \texttt{0400};
\item \texttt{0400 + 1200};
\item \texttt{0400 + IEG}.
\end{enumerate}

Three live trial families are proposed:
\begin{enumerate}[leftmargin=*]
\item no-touch routine pressure;
\item high-salience social pressure;
\item low-salience social pressure.
\end{enumerate}

This matrix is designed to test whether the simulator-side distinction survives
on real ERS-7 hardware and whether cue salience and multi-cycle persistence are
observable outside the host-side model.

\section{Contributions}
This work makes five current contributions:
\begin{enumerate}[leftmargin=*]
\item a verified recovery and preservation workflow for a repaired Sony-formatted
ERS-7 MIND 2 stick;
\item a reproducible behavior-forensics workflow over preserved persistent state;
\item a host-side simulator that exposes candidate behavioral levers from
persistent files;
\item a structural and dynamic characterization of specimen-side
\texttt{IEG.CFG};
\item a live-hardware matrix that turns current simulator conclusions into
falsifiable bench trials.
\end{enumerate}

\section{Limitations}
Several limits are important:
\begin{itemize}[leftmargin=*]
\item the simulator is heuristic and not a decoded Sony runtime;
\item current results are anchored in a small number of preserved reference
trees and one especially valuable repaired specimen;
\item the strongest findings are still host-side until repeated live-hardware
tests confirm or refute them;
\item some dynamic distinctions may reflect model choices rather than robot
truth and must therefore be treated as hypotheses.
\item current evidence supports live WLAN services and bounded content channels,
but not a general claim of arbitrary in-place retail stick reprogramming over
WLAN.
\end{itemize}

\section{Next Steps}
The next research steps are straightforward:
\begin{enumerate}[leftmargin=*]
\item execute the prepared four-stick live-hardware matrix;
\item log repeated observations with the current observation template;
\item compare live results to the current dynamic verdict tables;
\item refine the model only where live evidence demands it;
\item extend the paper with real-hardware results, figures, and stronger
bibliographic grounding in Sony documentation and robot behavior literature.
\end{enumerate}

\section{Conclusion}
The present work establishes that AIBO ERS-7 MIND 2 behavior can be approached
as a forensics problem over preserved filesystem state. A repaired and
preserved working stick can serve as a behavioral specimen. Host-side modeling
can isolate candidate persistent-state levers. Most importantly, current
evidence suggests that persistent behavior is not monolithic: shutdown
resistance, socially-cued persistence, and routine-preserving persistence are
separable enough to support structured experimentation. That is a meaningful
step toward a publishable science of legacy robotic behavior preservation and
reconstruction.

\end{document}
